\documentclass[12pt,a4paper]{article}

\usepackage[utf8]{inputenc}
\usepackage[T2A]{fontenc}
\usepackage[russian]{babel}
\usepackage{amsmath,amssymb,amsthm}
\usepackage{array}
\usepackage{booktabs}
\usepackage{enumitem}
\usepackage{listings}
\usepackage{xcolor}
\usepackage{geometry}

\geometry{margin=2.2cm}

\newtheorem{definition}{Определение}[section]
\newtheorem{theorem}{Теорема}[section]
\newtheorem{lemma}{Лемма}[section]
\newtheorem{proposition}{Утверждение}[section]
\newtheorem{example}{Пример}[section]

\theoremstyle{remark}
\newtheorem{remark}{Замечание}[section]

\lstdefinestyle{sqlstyle}{
    language=SQL,
    basicstyle=\ttfamily\small,
    keywordstyle=\color{blue},
    commentstyle=\color{gray},
    stringstyle=\color{red!60!black},
    breaklines=true,
    frame=single
}

\title{Билет №73 \\ \small CASE-средства и их использование при проектировании базы данных (БД), нотация IDEF. (СпБГУ)}
\author{Сериков Владислав Александрович}
\date{16 мая 2026}

\begin{document}

\maketitle
\tableofcontents
\newpage

\section{CASE-средства и их использование при проектировании баз данных. Нотации семейства IDEF}

\subsection{Основные понятия}

\textbf{CASE-средства} --- это программные средства, поддерживающие процессы анализа, проектирования, разработки, документирования, сопровождения и реинжиниринга информационных систем. Аббревиатура CASE обычно раскрывается как \textit{Computer-Aided Software Engineering}, то есть автоматизированная поддержка разработки программного обеспечения.

В контексте проектирования баз данных CASE-средства применяются для построения моделей данных, проверки их корректности, генерации схем БД, документирования структуры данных, сопровождения изменений и обратного проектирования уже существующих баз данных.

\textbf{База данных} --- это организованная совокупность взаимосвязанных данных, предназначенная для долговременного хранения, поиска, изменения и совместного использования.

\textbf{Система управления базами данных} --- программный комплекс, обеспечивающий создание, хранение, модификацию и защиту данных, а также доступ к ним со стороны прикладных программ и пользователей.

\textbf{Проектирование базы данных} --- процесс построения структуры данных, ограничений целостности и способов доступа к данным, удовлетворяющих требованиям предметной области и будущей информационной системы.

Обычно выделяют следующие уровни проектирования БД:

\begin{enumerate}
    \item \textbf{Инфологическое проектирование} --- описание предметной области независимо от конкретной СУБД. Результат: концептуальная модель данных.
    \item \textbf{Даталогическое проектирование} --- преобразование концептуальной модели в логическую модель, например в реляционную схему.
    \item \textbf{Физическое проектирование} --- выбор конкретных структур хранения, индексов, табличных пространств, параметров СУБД.
\end{enumerate}

CASE-средства особенно полезны на инфологическом и даталогическом этапах, поскольку позволяют формализовать знания о предметной области и уменьшить количество ошибок при переходе от требований к реализации.

\subsection{Классификация CASE-средств}

CASE-средства можно классифицировать по нескольким признакам.

\subsubsection{По этапам жизненного цикла}

\begin{enumerate}
    \item \textbf{Upper CASE} --- средства верхнего уровня, поддерживающие анализ предметной области, моделирование бизнес-процессов, построение концептуальных моделей.
    \item \textbf{Middle CASE} --- средства проектирования архитектуры, логических моделей данных, интерфейсов, спецификаций.
    \item \textbf{Lower CASE} --- средства генерации кода, схем БД, DDL-скриптов, тестов и документации.
    \item \textbf{Integrated CASE} --- интегрированные среды, поддерживающие несколько стадий жизненного цикла.
\end{enumerate}

\subsubsection{По назначению}

\begin{enumerate}
    \item средства функционального моделирования;
    \item средства информационного моделирования;
    \item средства объектно-ориентированного моделирования;
    \item средства проектирования реляционных баз данных;
    \item средства обратного проектирования;
    \item средства генерации документации;
    \item средства контроля версий и управления изменениями моделей.
\end{enumerate}

\subsubsection{По степени автоматизации}

\begin{enumerate}
    \item \textbf{Рисующие средства} --- позволяют строить диаграммы, но почти не проверяют их семантику.
    \item \textbf{Моделирующие средства} --- хранят модель в репозитории и проверяют ее внутреннюю согласованность.
    \item \textbf{Генерирующие средства} --- способны порождать SQL-код, DDL-скрипты, документацию, фрагменты программного кода.
    \item \textbf{Реинжиниринговые средства} --- восстанавливают модели по существующим схемам БД или программному коду.
\end{enumerate}

\subsection{Основные функции CASE-средств при проектировании БД}

CASE-средства для проектирования баз данных обычно поддерживают следующие функции:

\begin{enumerate}
    \item \textbf{Создание концептуальной модели данных}.

    На этом этапе описываются сущности предметной области, их атрибуты и связи между ними. Например, для университетской системы можно выделить сущности:
    \[
        \text{Студент}, \quad \text{Группа}, \quad \text{Дисциплина}, \quad \text{Экзамен}.
    \]

    \item \textbf{Создание логической модели данных}.

    Концептуальная модель преобразуется в логическую, обычно реляционную. Сущности становятся отношениями, атрибуты --- столбцами, связи --- внешними ключами.

    \item \textbf{Нормализация схемы}.

    CASE-средство может помогать выявлять избыточность данных, повторяющиеся группы, нарушения функциональных зависимостей.

    \item \textbf{Проверка ограничений целостности}.

    Средство проверяет наличие первичных ключей, корректность внешних ключей, обязательность атрибутов, типы связей, кардинальности.

    \item \textbf{Генерация DDL-скриптов}.

    По модели автоматически создаются SQL-команды:

\begin{verbatim}
CREATE TABLE Student (
    student_id INTEGER PRIMARY KEY,
    full_name  VARCHAR(200) NOT NULL,
    group_id   INTEGER NOT NULL,
    FOREIGN KEY (group_id) REFERENCES StudyGroup(group_id)
);
\end{verbatim}

    \item \textbf{Обратное проектирование}.

    CASE-средство может подключиться к существующей БД, считать структуру таблиц, ключей и связей, после чего построить диаграмму модели данных.

    \item \textbf{Синхронизация модели и базы данных}.

    При изменении модели можно сформировать миграционные скрипты, изменяющие физическую схему БД.

    \item \textbf{Документирование}.

    По модели формируются отчеты: список таблиц, атрибутов, ключей, ограничений, связей, комментариев.

    \item \textbf{Командная работа}.

    Современные CASE-средства могут поддерживать репозитории моделей, контроль версий, сравнение моделей и разрешение конфликтов.
\end{enumerate}

\subsection{Преимущества использования CASE-средств}

Использование CASE-средств при проектировании баз данных дает следующие преимущества:

\begin{enumerate}
    \item повышение наглядности проекта;
    \item снижение количества ошибок в схеме БД;
    \item автоматизация рутинных операций;
    \item ускорение перехода от модели к реализации;
    \item упрощение сопровождения БД;
    \item возможность коллективной разработки;
    \item единое хранилище проектной информации;
    \item автоматическое формирование документации;
    \item поддержка стандартных нотаций моделирования.
\end{enumerate}

\subsection{Недостатки и ограничения CASE-средств}

Несмотря на преимущества, CASE-средства не устраняют необходимость квалифицированного проектирования. Их ограничения:

\begin{enumerate}
    \item средство не заменяет анализ предметной области;
    \item ошибки аналитика могут быть формально корректно перенесены в модель;
    \item возможна зависимость от конкретного инструмента и его формата хранения моделей;
    \item сложные модели требуют дисциплины именования и документирования;
    \item автоматическая генерация схемы не всегда дает оптимальную физическую структуру;
    \item требуется обучение команды выбранной методологии и инструменту.
\end{enumerate}

\subsection{Типовой процесс проектирования БД с использованием CASE-средства}

Проектирование базы данных с помощью CASE-средства можно представить как алгоритм.

\subsubsection*{Алгоритм проектирования БД с использованием CASE-средства}

\begin{enumerate}
    \item \textbf{Сбор требований}.

    Определяются цели системы, пользователи, бизнес-процессы, документы, отчеты, правила предметной области.

    \item \textbf{Выделение объектов предметной области}.

    Из требований извлекаются основные сущности. Например:
    \[
        \text{Клиент}, \quad \text{Заказ}, \quad \text{Товар}, \quad \text{Платеж}.
    \]

    \item \textbf{Определение атрибутов сущностей}.

    Для каждой сущности задаются свойства:
    \[
        \text{Клиент}(\text{id}, \text{ФИО}, \text{телефон}, \text{email}).
    \]

    \item \textbf{Выбор идентификаторов}.

    Для каждой сущности выбирается первичный ключ. Например:
    \[
        \text{client\_id}
    \]
    для сущности \texttt{Client}.

    \item \textbf{Определение связей}.

    Устанавливаются связи между сущностями. Например:
    \[
        \text{Клиент} \; 1 : N \; \text{Заказ}.
    \]

    Это означает, что один клиент может иметь много заказов, но каждый заказ принадлежит одному клиенту.

    \item \textbf{Уточнение кардинальностей и обязательности}.

    Определяется, является ли участие сущности в связи обязательным. Например, заказ не может существовать без клиента, значит связь со стороны заказа обязательна.

    \item \textbf{Построение концептуальной модели}.

    В CASE-средстве строится ER-, UML- или IDEF1X-диаграмма.

    \item \textbf{Преобразование в логическую модель}.

    Сущности преобразуются в таблицы, атрибуты --- в столбцы, связи --- во внешние ключи.

    \item \textbf{Проверка модели}.

    Проверяются ключи, типы данных, связи, циклы, избыточные атрибуты, нарушения целостности.

    \item \textbf{Нормализация}.

    Схема приводится, как правило, по крайней мере к третьей нормальной форме, если нет специальных причин для денормализации.

    \item \textbf{Физическое проектирование}.

    Выбираются типы данных конкретной СУБД, индексы, ограничения, секционирование, параметры хранения.

    \item \textbf{Генерация SQL-скриптов}.

    CASE-средство формирует DDL-код для создания объектов БД.

    \item \textbf{Развертывание и тестирование}.

    Скрипты выполняются в тестовой БД, проверяются ограничения, запросы и производительность.

    \item \textbf{Документирование}.

    Генерируется описание схемы, диаграммы, словарь данных.

    \item \textbf{Сопровождение}.

    При изменении требований корректируется модель, а затем выполняется синхронизация с БД.
\end{enumerate}

\subsection{Минимальный пример проектирования БД в CASE-стиле}

Рассмотрим предметную область: студент записывается на учебные курсы.

\subsubsection*{Шаг 1. Выделение сущностей}

\[
    \text{Student}, \quad \text{Course}, \quad \text{Enrollment}.
\]

\subsubsection*{Шаг 2. Определение атрибутов}

\[
\begin{aligned}
    \text{Student} &= (\text{student\_id}, \text{full\_name}, \text{email}),\\
    \text{Course} &= (\text{course\_id}, \text{title}, \text{credits}),\\
    \text{Enrollment} &= (\text{student\_id}, \text{course\_id}, \text{enroll\_date}, \text{grade}).
\end{aligned}
\]

\subsubsection*{Шаг 3. Определение связей}

Один студент может быть записан на несколько курсов, и один курс может иметь много студентов. Следовательно, между \texttt{Student} и \texttt{Course} существует связь многие-ко-многим:

\[
    \text{Student} \; M : N \; \text{Course}.
\]

В реляционной модели такая связь реализуется через ассоциативную таблицу:

\[
    \text{Enrollment}.
\]

\subsubsection*{Шаг 4. Логическая схема}

\[
\begin{aligned}
    \text{Student}(&\underline{\text{student\_id}}, \text{full\_name}, \text{email}),\\
    \text{Course}(&\underline{\text{course\_id}}, \text{title}, \text{credits}),\\
    \text{Enrollment}(&\underline{\text{student\_id}, \text{course\_id}}, \text{enroll\_date}, \text{grade}).
\end{aligned}
\]

Здесь подчеркивание обозначает первичный ключ. В таблице \texttt{Enrollment} первичный ключ составной.

\subsubsection*{Шаг 5. DDL-код}

\begin{verbatim}
CREATE TABLE Student (
    student_id INTEGER PRIMARY KEY,
    full_name  VARCHAR(200) NOT NULL,
    email      VARCHAR(200) UNIQUE
);

CREATE TABLE Course (
    course_id INTEGER PRIMARY KEY,
    title     VARCHAR(200) NOT NULL,
    credits   INTEGER NOT NULL CHECK (credits > 0)
);

CREATE TABLE Enrollment (
    student_id  INTEGER NOT NULL,
    course_id   INTEGER NOT NULL,
    enroll_date DATE NOT NULL,
    grade       INTEGER CHECK (grade BETWEEN 2 AND 5),
    PRIMARY KEY (student_id, course_id),
    FOREIGN KEY (student_id) REFERENCES Student(student_id),
    FOREIGN KEY (course_id) REFERENCES Course(course_id)
);
\end{verbatim}

\subsection{Основные нотации моделирования данных}

При проектировании БД CASE-средства могут поддерживать разные нотации:

\begin{enumerate}
    \item ER-нотация Чена;
    \item нотация Crow's Foot;
    \item IDEF1X;
    \item UML Class Diagram;
    \item Barker notation;
    \item собственные нотации конкретных CASE-средств.
\end{enumerate}

Для проектирования реляционных баз данных особенно важна нотация \textbf{IDEF1X}, поскольку она ориентирована на информационное моделирование и хорошо согласуется с реляционной моделью данных.

\section{Семейство методологий IDEF}

\subsection{Общая характеристика IDEF}

\textbf{IDEF} --- семейство методологий моделирования, предназначенных для описания функций, процессов, информации, данных и поведения систем.

Наиболее известны:

\begin{enumerate}
    \item \textbf{IDEF0} --- функциональное моделирование;
    \item \textbf{IDEF1} --- информационное моделирование;
    \item \textbf{IDEF1X} --- моделирование данных;
    \item \textbf{IDEF3} --- описание процессов и сценариев;
    \item \textbf{IDEF4} --- объектно-ориентированное проектирование;
    \item \textbf{IDEF5} --- онтологическое моделирование.
\end{enumerate}

При проектировании баз данных наиболее важны IDEF0 и IDEF1X:

\begin{itemize}
    \item IDEF0 помогает понять функции и информационные потоки предметной области;
    \item IDEF1X используется для построения модели данных.
\end{itemize}

\section{Нотация IDEF0}

\subsection{Назначение IDEF0}

\textbf{IDEF0} --- методология функционального моделирования, применяемая для описания деятельности организации или системы. Она позволяет представить систему как совокупность взаимосвязанных функций.

IDEF0 полезна при проектировании БД, потому что позволяет выявить:

\begin{enumerate}
    \item какие процессы существуют в предметной области;
    \item какие данные используются процессами;
    \item какие документы и объекты поступают на вход;
    \item какие результаты формируются на выходе;
    \item какие правила управляют процессами;
    \item какие исполнители и средства обеспечивают выполнение функций.
\end{enumerate}

\subsection{Базовая конструкция IDEF0}

Основным элементом IDEF0 является \textbf{функциональный блок}. Он изображается прямоугольником, а стрелки вокруг него имеют строго заданную семантику.

\[
\begin{array}{c}
    \text{Управление} \\
    \downarrow \\
    \boxed{\text{Функция}} \\
    \uparrow \\
    \text{Механизм}
\end{array}
\]

Формально функцию можно описать как преобразование:

\[
    Y = F(X, C, M),
\]

где:

\begin{itemize}
    \item $X$ --- входы;
    \item $C$ --- управления;
    \item $M$ --- механизмы;
    \item $Y$ --- выходы.
\end{itemize}

В IDEF0 используется аббревиатура \textbf{ICOM}:

\[
    ICOM = Input, Control, Output, Mechanism.
\]

\subsection{Типы стрелок IDEF0}

\begin{enumerate}
    \item \textbf{Input} --- вход.

    Данные или материальные объекты, преобразуемые функцией.

    \item \textbf{Control} --- управление.

    Правила, ограничения, стандарты, инструкции, влияющие на выполнение функции, но не преобразуемые ею.

    \item \textbf{Output} --- выход.

    Результат выполнения функции.

    \item \textbf{Mechanism} --- механизм.

    Исполнители, технические средства, программные системы, ресурсы.
\end{enumerate}

\subsection{Пример IDEF0-модели для проектирования БД}

Рассмотрим процесс \textbf{оформления заказа}.

\[
\begin{array}{lll}
    \text{Входы:} & \text{заявка клиента, данные о товарах};\\
    \text{Управления:} & \text{правила скидок, договор, регламент продаж};\\
    \text{Выходы:} & \text{заказ, счет, резерв товара};\\
    \text{Механизмы:} & \text{менеджер, информационная система, складская система}.
\end{array}
\]

Из такой модели можно извлечь кандидаты в сущности БД:

\[
    \text{Client}, \quad \text{Order}, \quad \text{Product}, \quad \text{Invoice}, \quad \text{Warehouse}.
\]

\subsection{Декомпозиция в IDEF0}

IDEF0 поддерживает иерархическую декомпозицию функций. Верхний уровень описывает систему в целом, а нижние уровни раскрывают детали.

Например:

\[
    A0 = \text{Управлять продажами}.
\]

Декомпозиция:

\[
\begin{aligned}
    A1 &= \text{Принять заказ},\\
    A2 &= \text{Проверить наличие товара},\\
    A3 &= \text{Сформировать счет},\\
    A4 &= \text{Отгрузить товар}.
\end{aligned}
\]

Для проектирования БД это важно, поскольку на каждом уровне декомпозиции выявляются новые данные и связи.

\section{Нотация IDEF1X}

\subsection{Назначение IDEF1X}

\textbf{IDEF1X} --- нотация информационного моделирования, предназначенная для описания структуры данных предметной области. Она особенно удобна для проектирования реляционных баз данных.

IDEF1X используется для:

\begin{enumerate}
    \item выявления сущностей предметной области;
    \item описания атрибутов сущностей;
    \item задания идентификаторов;
    \item описания связей между сущностями;
    \item отображения обязательности и кардинальности связей;
    \item подготовки модели к преобразованию в реляционную схему.
\end{enumerate}

\subsection{Сущность в IDEF1X}

\textbf{Сущность} --- множество объектов предметной области одного типа, о которых необходимо хранить информацию.

Примеры сущностей:

\[
    \text{Студент}, \quad \text{Преподаватель}, \quad \text{Кафедра}, \quad \text{Дисциплина}.
\]

В IDEF1X сущность изображается прямоугольником, внутри которого указываются имя сущности и атрибуты.

Пример:

\[
\begin{array}{|l|}
\hline
\text{STUDENT}\\
\hline
\text{student\_id}\\
\hline
\text{full\_name}\\
\text{birth\_date}\\
\text{email}\\
\hline
\end{array}
\]

\subsection{Атрибут}

\textbf{Атрибут} --- именованное свойство сущности.

Например, для сущности \texttt{Student} атрибутами могут быть:

\[
    \text{student\_id}, \quad \text{full\_name}, \quad \text{birth\_date}, \quad \text{email}.
\]

Атрибуты бывают:

\begin{enumerate}
    \item \textbf{ключевые} --- входят в идентификатор сущности;
    \item \textbf{неключевые} --- описывают свойства сущности;
    \item \textbf{обязательные} --- не допускают отсутствующего значения;
    \item \textbf{необязательные} --- допускают отсутствие значения.
\end{enumerate}

\subsection{Идентификатор}

\textbf{Идентификатор сущности} --- один или несколько атрибутов, значения которых однозначно определяют экземпляр сущности.

Если идентификатор состоит из одного атрибута, он называется простым:

\[
    \text{student\_id}.
\]

Если идентификатор состоит из нескольких атрибутов, он называется составным:

\[
    (\text{student\_id}, \text{course\_id}).
\]

В реляционной схеме идентификатор обычно становится первичным ключом.

\subsection{Связи в IDEF1X}

\textbf{Связь} --- ассоциация между сущностями.

В IDEF1X обычно различают:

\begin{enumerate}
    \item идентифицирующие связи;
    \item неидентифицирующие связи;
    \item обязательные связи;
    \item необязательные связи;
    \item связи один-к-одному;
    \item связи один-ко-многим;
    \item связи многие-ко-многим.
\end{enumerate}

\subsection{Родительская и дочерняя сущность}

В связи IDEF1X выделяют:

\begin{itemize}
    \item \textbf{родительскую сущность} --- сущность, чей ключ передается в другую сущность;
    \item \textbf{дочернюю сущность} --- сущность, которая получает внешний ключ.
\end{itemize}

Например, если один факультет имеет много кафедр:

\[
    \text{Faculty} \; 1 : N \; \text{Department},
\]

то \texttt{Faculty} является родительской сущностью, а \texttt{Department} --- дочерней.

\subsection{Идентифицирующая связь}

\textbf{Идентифицирующая связь} --- связь, при которой ключ родительской сущности входит в состав первичного ключа дочерней сущности.

Например:

\[
\begin{aligned}
    \text{Order}(&\underline{\text{order\_id}}, \text{order\_date}),\\
    \text{OrderItem}(&\underline{\text{order\_id}, \text{line\_no}}, \text{product\_id}, \text{quantity}).
\end{aligned}
\]

Здесь \texttt{order\_id} входит в первичный ключ таблицы \texttt{OrderItem}. Поэтому связь между \texttt{Order} и \texttt{OrderItem} является идентифицирующей.

Смысл: строка заказа не имеет самостоятельной идентичности вне заказа.

\subsection{Неидентифицирующая связь}

\textbf{Неидентифицирующая связь} --- связь, при которой ключ родительской сущности мигрирует в дочернюю сущность как внешний ключ, но не входит в ее первичный ключ.

Например:

\[
\begin{aligned}
    \text{Department}(&\underline{\text{department\_id}}, \text{name}),\\
    \text{Employee}(&\underline{\text{employee\_id}}, \text{full\_name}, \text{department\_id}).
\end{aligned}
\]

Здесь \texttt{department\_id} является внешним ключом в \texttt{Employee}, но не входит в первичный ключ сотрудника. Следовательно, связь неидентифицирующая.

\subsection{Зависимые и независимые сущности}

В IDEF1X различают:

\begin{enumerate}
    \item \textbf{независимую сущность};
    \item \textbf{зависимую сущность}.
\end{enumerate}

\textbf{Независимая сущность} имеет собственный идентификатор, не зависящий от других сущностей.

Пример:

\[
    \text{Student}(\underline{\text{student\_id}}, \text{full\_name}).
\]

\textbf{Зависимая сущность} идентифицируется с использованием ключа другой сущности.

Пример:

\[
    \text{OrderItem}(\underline{\text{order\_id}, \text{line\_no}}, \text{quantity}).
\]

Здесь \texttt{OrderItem} зависит от \texttt{Order}.

\subsection{Миграция ключей}

\textbf{Миграция ключа} --- перенос первичного ключа родительской сущности в дочернюю сущность в виде внешнего ключа.

Например:

\[
\begin{aligned}
    \text{Client}(&\underline{\text{client\_id}}, \text{name}),\\
    \text{Order}(&\underline{\text{order\_id}}, \text{order\_date}, \text{client\_id}).
\end{aligned}
\]

Ключ \texttt{client\_id} мигрирует из \texttt{Client} в \texttt{Order}.

В SQL это выражается ограничением:

\begin{verbatim}
FOREIGN KEY (client_id) REFERENCES Client(client_id)
\end{verbatim}

\subsection{Кардинальность связей}

Кардинальность показывает, сколько экземпляров одной сущности может быть связано с экземпляром другой сущности.

Основные типы:

\begin{enumerate}
    \item \textbf{один-к-одному}:
    \[
        1 : 1;
    \]

    \item \textbf{один-ко-многим}:
    \[
        1 : N;
    \]

    \item \textbf{многие-ко-многим}:
    \[
        M : N.
    \]
\end{enumerate}

В реляционной модели связь многие-ко-многим непосредственно не реализуется. Она преобразуется в отдельную ассоциативную сущность.

\subsection{Преобразование связи многие-ко-многим}

Пусть имеется связь:

\[
    \text{Student} \; M : N \; \text{Course}.
\]

Она преобразуется в две связи один-ко-многим через ассоциативную сущность:

\[
    \text{Student} \; 1 : N \; \text{Enrollment} \; N : 1 \; \text{Course}.
\]

Логическая схема:

\[
\begin{aligned}
    \text{Student}(&\underline{\text{student\_id}}, \text{full\_name}),\\
    \text{Course}(&\underline{\text{course\_id}}, \text{title}),\\
    \text{Enrollment}(&\underline{\text{student\_id}, \text{course\_id}}, \text{enroll\_date}).
\end{aligned}
\]

\section{Алгоритмы преобразования моделей IDEF1X в реляционную схему}

\subsection{Алгоритм преобразования сущностей в таблицы}

\subsubsection*{Вход}

IDEF1X-модель, содержащая множество сущностей:

\[
    E = \{E_1, E_2, \ldots, E_n\}.
\]

\subsubsection*{Выход}

Реляционная схема:

\[
    R = \{R_1, R_2, \ldots, R_n\}.
\]

\subsubsection*{Шаги алгоритма}

\begin{enumerate}
    \item Для каждой сущности $E_i$ создать отношение $R_i$.
    \item Для каждого атрибута сущности создать столбец соответствующего отношения.
    \item Атрибуты идентификатора сущности назначить первичным ключом.
    \item Для обязательных атрибутов задать ограничение \texttt{NOT NULL}.
    \item Для доменов атрибутов выбрать типы данных конкретной СУБД.
    \item Для уникальных альтернативных идентификаторов задать ограничение \texttt{UNIQUE}.
\end{enumerate}

\subsubsection*{Минимальный пример}

IDEF1X-сущность:

\[
    \text{Book}(\underline{\text{book\_id}}, \text{title}, \text{isbn}).
\]

Реляционная таблица:

\begin{verbatim}
CREATE TABLE Book (
    book_id INTEGER PRIMARY KEY,
    title   VARCHAR(300) NOT NULL,
    isbn    VARCHAR(20) UNIQUE
);
\end{verbatim}

\subsection{Алгоритм преобразования связи один-ко-многим}

\subsubsection*{Вход}

Связь:

\[
    A \; 1 : N \; B,
\]

где $A$ --- родительская сущность, $B$ --- дочерняя сущность.

\subsubsection*{Выход}

Реляционные таблицы с внешним ключом.

\subsubsection*{Шаги алгоритма}

\begin{enumerate}
    \item Создать таблицу для сущности $A$.
    \item Создать таблицу для сущности $B$.
    \item Первичный ключ таблицы $A$ добавить в таблицу $B$ как внешний ключ.
    \item Если связь обязательна со стороны $B$, внешний ключ сделать \texttt{NOT NULL}.
    \item Если связь идентифицирующая, внешний ключ включить в первичный ключ таблицы $B$.
    \item Если связь неидентифицирующая, внешний ключ не включать в первичный ключ таблицы $B$.
\end{enumerate}

\subsubsection*{Пример неидентифицирующей связи}

\[
    \text{Department} \; 1 : N \; \text{Employee}.
\]

\begin{verbatim}
CREATE TABLE Department (
    department_id INTEGER PRIMARY KEY,
    name          VARCHAR(200) NOT NULL
);

CREATE TABLE Employee (
    employee_id   INTEGER PRIMARY KEY,
    full_name     VARCHAR(200) NOT NULL,
    department_id INTEGER NOT NULL,
    FOREIGN KEY (department_id)
        REFERENCES Department(department_id)
);
\end{verbatim}

\subsubsection*{Пример идентифицирующей связи}

\[
    \text{Order} \; 1 : N \; \text{OrderItem}.
\]

\begin{verbatim}
CREATE TABLE Orders (
    order_id   INTEGER PRIMARY KEY,
    order_date DATE NOT NULL
);

CREATE TABLE OrderItem (
    order_id  INTEGER NOT NULL,
    line_no   INTEGER NOT NULL,
    product   VARCHAR(200) NOT NULL,
    quantity  INTEGER NOT NULL,
    PRIMARY KEY (order_id, line_no),
    FOREIGN KEY (order_id)
        REFERENCES Orders(order_id)
);
\end{verbatim}

\subsection{Алгоритм преобразования связи многие-ко-многим}

\subsubsection*{Вход}

Связь:

\[
    A \; M : N \; B.
\]

\subsubsection*{Выход}

Три отношения: два основных и одно ассоциативное.

\subsubsection*{Шаги алгоритма}

\begin{enumerate}
    \item Создать таблицу для сущности $A$.
    \item Создать таблицу для сущности $B$.
    \item Создать новую ассоциативную таблицу $AB$.
    \item Добавить в $AB$ первичный ключ $A$ как внешний ключ.
    \item Добавить в $AB$ первичный ключ $B$ как внешний ключ.
    \item Составить первичный ключ $AB$ из внешних ключей $A$ и $B$ либо ввести отдельный суррогатный ключ.
    \item Добавить атрибуты, относящиеся к самой связи.
\end{enumerate}

\subsubsection*{Пример}

\[
    \text{Author} \; M : N \; \text{Book}.
\]

\begin{verbatim}
CREATE TABLE Author (
    author_id INTEGER PRIMARY KEY,
    full_name VARCHAR(200) NOT NULL
);

CREATE TABLE Book (
    book_id INTEGER PRIMARY KEY,
    title   VARCHAR(300) NOT NULL
);

CREATE TABLE AuthorBook (
    author_id INTEGER NOT NULL,
    book_id   INTEGER NOT NULL,
    role      VARCHAR(100),
    PRIMARY KEY (author_id, book_id),
    FOREIGN KEY (author_id) REFERENCES Author(author_id),
    FOREIGN KEY (book_id) REFERENCES Book(book_id)
);
\end{verbatim}

\section{Место нормализации в CASE-проектировании}

Нормализация --- процесс организации отношений в реляционной БД с целью уменьшения избыточности и предотвращения аномалий вставки, обновления и удаления.

В CASE-проектировании нормализация обычно выполняется после построения логической модели.

\subsection{Первая нормальная форма}

Отношение находится в первой нормальной форме, если все значения его атрибутов атомарны.

Нарушение:

\[
    \text{Student}(\text{id}, \text{name}, \text{phones}).
\]

Если в атрибуте \texttt{phones} хранится список телефонов, то значение неатомарно.

Исправление:

\[
\begin{aligned}
    \text{Student}(&\underline{\text{id}}, \text{name}),\\
    \text{StudentPhone}(&\underline{\text{id}, \text{phone}}).
\end{aligned}
\]

\subsection{Вторая нормальная форма}

Отношение находится во второй нормальной форме, если оно находится в первой нормальной форме и каждый неключевой атрибут полностью функционально зависит от всего первичного ключа, а не от его части.

Пример нарушения:

\[
    \text{Enrollment}(\underline{\text{student\_id}, \text{course\_id}}, \text{student\_name}, \text{course\_title}, \text{grade}).
\]

Здесь:

\[
    \text{student\_id} \rightarrow \text{student\_name},
\]
\[
    \text{course\_id} \rightarrow \text{course\_title}.
\]

Следовательно, \texttt{student\_name} и \texttt{course\_title} зависят только от части составного ключа.

Исправление:

\[
\begin{aligned}
    \text{Student}(&\underline{\text{student\_id}}, \text{student\_name}),\\
    \text{Course}(&\underline{\text{course\_id}}, \text{course\_title}),\\
    \text{Enrollment}(&\underline{\text{student\_id}, \text{course\_id}}, \text{grade}).
\end{aligned}
\]

\subsection{Третья нормальная форма}

Отношение находится в третьей нормальной форме, если оно находится во второй нормальной форме и неключевые атрибуты не зависят транзитивно от первичного ключа.

Пример нарушения:

\[
    \text{Employee}(\underline{\text{employee\_id}}, \text{name}, \text{department\_id}, \text{department\_name}).
\]

Имеются зависимости:

\[
    \text{employee\_id} \rightarrow \text{department\_id},
\]
\[
    \text{department\_id} \rightarrow \text{department\_name}.
\]

Следовательно:

\[
    \text{employee\_id} \rightarrow \text{department\_name}
\]

транзитивно.

Исправление:

\[
\begin{aligned}
    \text{Employee}(&\underline{\text{employee\_id}}, \text{name}, \text{department\_id}),\\
    \text{Department}(&\underline{\text{department\_id}}, \text{department\_name}).
\end{aligned}
\]

\section{Связь CASE-средств, IDEF0 и IDEF1X}

При проектировании информационной системы удобно использовать несколько моделей одновременно.

\subsection{Общая схема}

\[
\begin{array}{ccccc}
\text{Требования}
& \rightarrow &
\text{IDEF0-модель процессов}
& \rightarrow &
\text{IDEF1X-модель данных}
\\
& &
\downarrow
& &
\downarrow
\\
& &
\text{Информационные потоки}
& \rightarrow &
\text{Реляционная схема БД}
\end{array}
\]

\subsection{Роль IDEF0}

IDEF0 отвечает на вопросы:

\begin{enumerate}
    \item какие функции выполняет система;
    \item какие данные поступают на вход;
    \item какие данные формируются на выходе;
    \item какие правила управляют процессами;
    \item какие исполнители и средства участвуют в работе.
\end{enumerate}

\subsection{Роль IDEF1X}

IDEF1X отвечает на вопросы:

\begin{enumerate}
    \item какие сущности необходимо хранить;
    \item какие атрибуты они имеют;
    \item какие ключи идентифицируют экземпляры;
    \item какие связи существуют между сущностями;
    \item как эти связи будут реализованы в реляционной БД.
\end{enumerate}

\subsection{Роль CASE-средства}

CASE-средство обеспечивает:

\begin{enumerate}
    \item хранение моделей;
    \item визуальное редактирование диаграмм;
    \item проверку согласованности;
    \item генерацию SQL-кода;
    \item документирование;
    \item обратное проектирование;
    \item сопровождение изменений.
\end{enumerate}

\section{Пример сквозного проектирования}

Рассмотрим информационную систему библиотеки.

\subsection{IDEF0-уровень}

Основная функция:

\[
    A0 = \text{Обслуживать читателей библиотеки}.
\]

Входы:

\[
    \text{заявка на выдачу книги}, \quad \text{возвращаемая книга}.
\]

Управления:

\[
    \text{правила выдачи}, \quad \text{ограничения по срокам}, \quad \text{каталог библиотеки}.
\]

Выходы:

\[
    \text{выданная книга}, \quad \text{запись о выдаче}, \quad \text{уведомление о просрочке}.
\]

Механизмы:

\[
    \text{библиотекарь}, \quad \text{информационная система}, \quad \text{читательский билет}.
\]

\subsection{Выделение сущностей}

Из функциональной модели можно выделить сущности:

\[
    \text{Reader}, \quad \text{Book}, \quad \text{BookCopy}, \quad \text{Loan}, \quad \text{Librarian}.
\]

\subsection{IDEF1X-модель}

\[
\begin{aligned}
    \text{Reader}(&\underline{\text{reader\_id}}, \text{full\_name}, \text{phone}),\\
    \text{Book}(&\underline{\text{book\_id}}, \text{title}, \text{isbn}),\\
    \text{BookCopy}(&\underline{\text{copy\_id}}, \text{book\_id}, \text{inventory\_no}),\\
    \text{Librarian}(&\underline{\text{librarian\_id}}, \text{full\_name}),\\
    \text{Loan}(&\underline{\text{loan\_id}}, \text{reader\_id}, \text{copy\_id}, \text{librarian\_id}, \text{loan\_date}, \text{return\_date}).
\end{aligned}
\]

\subsection{Связи}

\[
    \text{Book} \; 1 : N \; \text{BookCopy}.
\]

Одна книга может иметь несколько физических экземпляров.

\[
    \text{Reader} \; 1 : N \; \text{Loan}.
\]

Один читатель может иметь много выдач.

\[
    \text{BookCopy} \; 1 : N \; \text{Loan}.
\]

Один экземпляр книги может участвовать во многих выдачах в разные моменты времени.

\[
    \text{Librarian} \; 1 : N \; \text{Loan}.
\]

Один библиотекарь может оформить много выдач.

\subsection{DDL-фрагмент}

\begin{verbatim}
CREATE TABLE Reader (
    reader_id INTEGER PRIMARY KEY,
    full_name VARCHAR(200) NOT NULL,
    phone     VARCHAR(30)
);

CREATE TABLE Book (
    book_id INTEGER PRIMARY KEY,
    title   VARCHAR(300) NOT NULL,
    isbn    VARCHAR(20) UNIQUE
);

CREATE TABLE BookCopy (
    copy_id      INTEGER PRIMARY KEY,
    book_id      INTEGER NOT NULL,
    inventory_no VARCHAR(50) NOT NULL UNIQUE,
    FOREIGN KEY (book_id) REFERENCES Book(book_id)
);

CREATE TABLE Librarian (
    librarian_id INTEGER PRIMARY KEY,
    full_name    VARCHAR(200) NOT NULL
);

CREATE TABLE Loan (
    loan_id      INTEGER PRIMARY KEY,
    reader_id    INTEGER NOT NULL,
    copy_id      INTEGER NOT NULL,
    librarian_id INTEGER NOT NULL,
    loan_date    DATE NOT NULL,
    return_date  DATE,
    FOREIGN KEY (reader_id) REFERENCES Reader(reader_id),
    FOREIGN KEY (copy_id) REFERENCES BookCopy(copy_id),
    FOREIGN KEY (librarian_id) REFERENCES Librarian(librarian_id)
);
\end{verbatim}

\section{Примеры CASE-средств для проектирования БД}

К распространенным CASE- и data modeling-средствам относятся:

\begin{enumerate}
    \item ERwin Data Modeler;
    \item SAP PowerDesigner;
    \item Oracle SQL Developer Data Modeler;
    \item IBM InfoSphere Data Architect;
    \item Enterprise Architect;
    \item Visual Paradigm;
    \item pgModeler;
    \item MySQL Workbench;
    \item DBeaver с ER-диаграммами;
    \item DataGrip с диаграммами зависимостей.
\end{enumerate}

Типичные возможности таких средств:

\begin{enumerate}
    \item построение ER-, UML- или IDEF1X-диаграмм;
    \item описание сущностей, таблиц, атрибутов, ключей;
    \item генерация DDL;
    \item reverse engineering существующей БД;
    \item сравнение модели и физической БД;
    \item генерация отчетов;
    \item поддержка нескольких СУБД;
    \item настройка типов данных для конкретной СУБД.
\end{enumerate}

\section{Критерии качества модели данных}

Хорошая модель данных должна обладать следующими свойствами:

\begin{enumerate}
    \item \textbf{полнота} --- отражает все существенные объекты предметной области;
    \item \textbf{непротиворечивость} --- не содержит логических конфликтов;
    \item \textbf{минимальная избыточность} --- не дублирует данные без необходимости;
    \item \textbf{нормализованность} --- устраняет основные аномалии данных;
    \item \textbf{расширяемость} --- допускает развитие при изменении требований;
    \item \textbf{реализуемость} --- может быть преобразована в схему конкретной СУБД;
    \item \textbf{понятность} --- читаема для аналитиков, разработчиков и администраторов;
    \item \textbf{трассируемость} --- элементы модели связаны с требованиями и бизнес-процессами.
\end{enumerate}

\section{Типичные ошибки при использовании CASE-средств}

\begin{enumerate}
    \item построение диаграммы без предварительного анализа предметной области;
    \item смешение концептуального, логического и физического уровней;
    \item использование технических идентификаторов без бизнес-смысла там, где нужны естественные ключи;
    \item отсутствие описаний сущностей и атрибутов;
    \item некорректное моделирование связей многие-ко-многим;
    \item чрезмерная денормализация на раннем этапе;
    \item игнорирование ограничений целостности;
    \item отсутствие единых правил именования;
    \item ручное изменение БД без синхронизации с моделью;
    \item использование CASE-средства только как редактора картинок.
\end{enumerate}

\section{Выводы}

CASE-средства являются важным инструментом проектирования баз данных. Они позволяют формализовать требования, построить концептуальную и логическую модели, проверить корректность структуры данных, сгенерировать SQL-код и документацию.

Методологии семейства IDEF полезны на разных этапах проектирования. IDEF0 применяется для функционального анализа предметной области и выявления информационных потоков. IDEF1X применяется для детального моделирования данных, сущностей, атрибутов, ключей и связей.

Наиболее эффективный подход состоит в последовательном переходе:

\[
    \text{требования}
    \rightarrow
    \text{функциональная модель}
    \rightarrow
    \text{информационная модель}
    \rightarrow
    \text{логическая схема}
    \rightarrow
    \text{физическая БД}.
\]

CASE-средство не заменяет аналитика и проектировщика, но существенно повышает качество, согласованность и сопровождаемость проекта базы данных.

\end{document}
